\documentclass[11pt]{article}
\usepackage[margin=1in]{geometry}
\usepackage{amsmath,amssymb,amsthm}
\usepackage{hyperref}
\usepackage{enumitem}

\newtheorem{theorem}{Theorem}
\newtheorem{lemma}[theorem]{Lemma}
\newtheorem{corollary}[theorem]{Corollary}
\theoremstyle{definition}
\newtheorem{definition}[theorem]{Definition}
\theoremstyle{remark}
\newtheorem{remark}[theorem]{Remark}

\usepackage{xurl}
\let\oldus\_
\renewcommand{\_}{\oldus\allowbreak}
\newcommand{\repdim}{\operatorname{repdim}}
\newcommand{\gldim}{\operatorname{gl.dim}}
\newcommand{\pd}{\operatorname{pd}}
\newcommand{\add}{\operatorname{add}}
\newcommand{\End}{\operatorname{End}}
\newcommand{\Hom}{\operatorname{Hom}}

\title{A disproof of Conjecture 00000004042\\[2pt]
\large The values of the representation dimension are not dense above $3$}
\date{}

\begin{document}
\sloppy
\maketitle

\section{The conjecture}

Conjecture 00000004042 reads (English version):

\begin{quote}
\emph{Definition:} An interval algebra is an algebra constructed from a poset with repdim = 3.
\emph{Conjecture:} The realizable values of repdim are dense above 3, and the class of algebras with
repdim = 3 is closed under derived equivalence.
\end{quote}

The Chinese version (in \texttt{conjectures/00000004042.md}) says, in transliteration,
``repdim de ke shixian zhi \emph{zai 3 yishang choumi}; repdim = 3 de daishu lei zai daochu dengjia xia
fengbi'', i.e.\ ``the realizable values of repdim are \emph{dense in [the region] above 3}; the class of
algebras with repdim $=3$ is closed under derived equivalence''.

The conjecture is the conjunction of two clauses:
\begin{enumerate}[label=(\arabic*)]
  \item the set $V$ of realizable values of $\repdim$ is dense above $3$;
  \item if $\repdim A = 3$ and $B$ is derived equivalent to $A$, then $\repdim B = 3$.
\end{enumerate}
The preliminary ``Definition'' only introduces the name ``interval algebra''. Neither clause mentions
interval algebras, so the definition plays no role.

\textbf{Result.} Clause (1) is false: $\repdim$ takes values in $\{0,1,2,\dots\}\cup\{\infty\}$ (and
$-\infty$ for the zero algebra under Mathlib's conventions), so the nonempty open interval
$(13/4, 15/4) \subset (3,\infty)$ contains no realizable value. Hence the conjunction (1)$\wedge$(2) is
false, whatever clause (2) means. The proof is formalized in Lean~4 with Mathlib (Section~\ref{sec:lean}).

\section{Reading}\label{sec:reading}

\paragraph{repdim.} ``repdim'' is Auslander's representation dimension, the standard meaning of this
abbreviation in the representation theory of algebras (Definition~\ref{def:repdim}). The algebras are
finite-dimensional algebras over a field, the setting of the definition quoted below from
Erdmann--Holm--Iyama--Schr\"oer.

\paragraph{``Dense above 3''.} ``Dense'' is a technical term with a standard meaning. Wikipedia,
\emph{Dense set}: ``a subset A of a topological space X is said to be dense in X if every point of X
either belongs to A or else is arbitrarily ``close'' to members of A'' and ``Formally, $A$ is dense in
$X$ if the smallest closed subset of $X$ containing $A$ is $X$ itself.'' The Chinese text uses the
construction ``zai $X$ choumi'' (``dense in $X$''), which names the ambient space $X$, here ``3 yishang''
(``above 3''). The realizable values are numbers, so the ambient space is the half-line of real numbers
above $3$. Clause (1) therefore says:
\[
  \text{(1)}\qquad V \cap (3,\infty) \text{ is dense in the topological space } (3,\infty) \subset \mathbb{R}.
\]
Equivalently, every nonempty open subset of $(3,\infty)$ contains a realizable value. This is exactly
what the Lean definition \texttt{DenseAbove3} states, using Mathlib's \texttt{Dense} in the subspace
\texttt{Set.Ioi 3} of $\mathbb{R}$.

Two remarks on this reading.
\begin{itemize}
  \item Density passes to open subspaces: if $V$ were dense in $[3,\infty)$, or in the extended
        half-line $(3,+\infty]$, then it would be dense in the open subspace $(3,\infty)$. So the
        statement refuted here is implied by these variants, and the refutation covers them too.
  \item Scope of this reading. A different, more charitable reading is conceivable: if the ambient space
        were taken to be the discrete set of integers $\ge 3$, density there would amount to ``every
        integer $\ge 3$ is a realizable value''. Neither language version names such an ambient space or
        replaces density by integer realizability, and the realizable values are compared with the real
        number $3$; we therefore take the standard real-topological meaning stated above as the literal
        one. We state explicitly that this submission does not address the integer-realizability
        reading (nor clause (2) on its own); it refutes the conjecture as written, under the standard
        meaning of ``dense''. Likewise ``dense above 3'' is not read as ``densely ordered'': the
        construction ``dense in $X$'' (``zai $X$ choumi'') is the topological notion relative to the named
        ambient space $X$.
\end{itemize}

\paragraph{Clause (2).} Following Wikipedia, \emph{Tilting theory}, ``A and B are derived equivalent
(i.e. the derived categories Db(A-mod) and Db(B-mod) are equivalent as triangulated categories)''.
Because clause (1) is false, the conjunction is false whatever clause (2) says. The Lean development
makes this explicit. It proves $\neg(\texttt{DenseAbove3}\wedge P)$ for \emph{every} proposition $P$
(\texttt{conjecture\_00000004042\_false\_of\_any\_clause2}). The main theorem is the instance of this
at a concrete formalization \texttt{ClosedUnderDerivedEquivalence} of clause (2), built on Mathlib's
derived categories (Section~\ref{sec:lean}). So the refutation does not depend on any choice made in
formalizing clause (2): any other formalization (bounded derived categories of finitely generated
modules, $k$-linear equivalences, \dots) is also an instance of the general theorem. The proof never
uses the content of clause (2), so no fact about derived equivalence is needed.

\section{Definitions and sources}

All quotations below are verbatim from the retrieved sources. Mathematical symbols are transcribed
into \LaTeX.

\begin{definition}[Projective and global dimension]\label{def:gldim}
The projective dimension $\pd X$ of a module $X$ is the minimal length of a projective resolution of
$X$, or $\infty$. Wikipedia, \emph{Resolution (algebra)}: ``The minimal length of a finite projective
resolution of a module $M$ is called its projective dimension'' and ``If $M$ does not admit a finite
projective resolution then the projective dimension is infinite.'' The (left) global dimension of a ring
$R$ is
\[
  \gldim R = \sup\{\pd X : X \text{ a left } R\text{-module}\}.
\]
Wikipedia, \emph{Global dimension}: the global dimension ``of a ring A denoted gl dim A, is a non-negative
integer or infinity which is a homological invariant of the ring. It is defined to be the supremum of
the set of projective dimensions of all A-modules.''
\end{definition}

In Lean we use Mathlib's \texttt{CategoryTheory.projectiveDimension X}, which takes values in
\texttt{WithBot $\mathbb{N}_\infty$} $= \{\bot\}\cup\{0,1,2,\dots\}\cup\{\infty\}$. Its docstring
reads: ``We also define the projective dimension in \texttt{WithBot $\mathbb{N}_\infty$} as
\texttt{projectiveDimension}, \texttt{projectiveDimension X = $\bot$} iff \texttt{X} is zero and behaves
as expected on non-negative values''. (``$X$ has projective dimension $< n$'' means that
$\mathrm{Ext}^i(X,-)=0$ for all $i \ge n$.) The value $\bot$ occurs only for the zero module.

\begin{definition}[Generator-cogenerator, representation dimension]\label{def:repdim}
Let $A$ be a finite-dimensional algebra over a field $k$, and let $D = \Hom_k(-,k)$. For an $A$-module
$M$, $\add M$ is the class of direct summands of finite direct sums $M^n$. A
\emph{generator-cogenerator} is a finitely generated $A$-module $M$ with $A \in \add M$ and
$D(A_A) \in \add M$ (equivalently $A \oplus D(A_A) \in \add M$). The \emph{representation dimension} is
\[
  \repdim A = \inf\{\gldim \End_A(M) : M \text{ a generator-cogenerator of } A\}.
\]
\end{definition}

Sources:
\begin{itemize}
  \item K.~Erdmann, T.~Holm, O.~Iyama, J.~Schr\"oer, \emph{Radical embeddings and representation
        dimension}, arXiv:math/0210362, Section~1 (\url{https://arxiv.org/html/math/0210362}):
        ``Throughout, let k be a field, and let A be an associative finite-dimensional k-algebra. By
        mod(A) we denote the category of finite-dimensional left A-modules. A generator-cogenerator of A
        is an A-module Z such that $A \oplus A^* \in \add(Z)$, where $A^* = D(A_A)$ with
        $D = \Hom_k(-,k)$ the usual duality. Recall that add(Z) is the subcategory of all A-modules which
        are isomorphic to direct summands of finite direct sums of Z. The representation dimension of A
        is defined as $\repdim(A) = \min\{\operatorname{gldim}(\End_A(Z)) \mid Z$ a generator-cogenerator
        of $A\}$.''
  \item W.~Hu, C.~C.~Xi, \emph{Derived equivalences and stable equivalences of Morita type, I},
        arXiv:0810.4761v2, Section~5 (\url{https://arxiv.org/html/0810.4761v2}): ``An A-module M is
        called a generator-cogenerator for A-mod if $A \oplus D(A) \in \add(M)$. The representation
        dimension of A, denoted by rep.dim(A), is defined to be
        $\mathrm{rep.dim}(A) := \inf\{\mathrm{gl.dim}(\End_A(M)) \mid M$ is a generator-cogenerator for
        A-mod$\}$. This notion was introduced by Auslander in [1] to measure homologically how far an
        algebra is from being representation-finite''.
  \item C.~M.~Ringel, \emph{On the representation dimension of artin algebras}, arXiv:1107.1861v1,
        abstract (\url{https://arxiv.org/html/1107.1861v1}): ``The representation dimension of an artin
        algebra as introduced by M. Auslander in his Queen Mary Notes is the minimal possible global
        dimension of the endomorphism ring of a generator-cogenerator.''
\end{itemize}
EHIS write $\min$, Hu--Xi write $\inf$. The proof below uses only that $\repdim A$ is an infimum of
global dimensions, so it applies to both formulations.

\begin{remark}[Conventions that do not affect the argument]\label{rem:conv}
(a) A finitely generated module over a finite-dimensional algebra is finite-dimensional over $k$, and
conversely, so ``finitely generated'' and EHIS's ``finite-dimensional'' agree.
(b) $\End_A(M)$ acts on $M$ on the left, and we use its left global dimension. $\End_A(M)$ is a
finite-dimensional $k$-algebra, hence Noetherian on both sides, and for such rings left and right global
dimensions coincide (Wikipedia, \emph{Global dimension}: ``if A is a Noetherian ring, both of these
dimensions turn out to be equal to weak global dimension''). In any case every global dimension, left or
right, lies in $\{0,1,2,\dots\}\cup\{\infty\}$, and that is all the proof uses.
(c) Universes. In Lean, the field $k$, the algebra $A$, the modules $M$, and the modules over
$\End_A(M)$ all live in one universe $u$. This loses nothing. Every finitely generated $A$-module is
isomorphic to a quotient of some $A^n$, which lives in universe $u$, so every isomorphism class of
generator-cogenerators is represented. The global dimension of a ring is determined by its cyclic
modules, which are quotients of the ring and so live in its universe (Wikipedia, \emph{Global
dimension}: the right global dimension ``can be alternatively defined as: the supremum of the set of
projective dimensions of all cyclic right A-modules'', and ``The left global dimension of A has
analogous characterizations''). All Lean statements are universe polymorphic, and the theorem is proved
in every universe $u$.
\end{remark}

\section{The proof}

\begin{lemma}[Integrality]\label{lem:int}
For every finite-dimensional algebra $A$ over a field $k$, $\repdim A \in \{-\infty\}\cup\{0,1,2,\dots\}\cup\{\infty\}$.
In particular, every finite real value of $\repdim$ is a non-negative integer.
\end{lemma}

\begin{proof}
Every projective dimension lies in $W := \{\bot\}\cup\{0,1,2,\dots\}\cup\{\infty\}$, a complete linear
order. $\gldim R$ is a supremum of projective dimensions, so $\gldim R \in W$. $\repdim A$ is an
infimum of global dimensions, so $\repdim A \in W$. Under the order embedding
$W \hookrightarrow [-\infty,+\infty]$ ($\bot\mapsto-\infty$, $n\mapsto n$, $\infty\mapsto+\infty$), the only
elements of $W$ that map to real numbers are the natural numbers.
\end{proof}

\begin{lemma}\label{lem:gap}
No integer lies in the open interval $(13/4, 15/4)$.
\end{lemma}

\begin{proof}
If $13/4 < m < 15/4$ then $3 < m < 4$, which is impossible for an integer $m$.
\end{proof}

\begin{theorem}\label{thm:main}
Clause (1) is false: the set $V$ of realizable values of $\repdim$ is not dense in $(3,\infty)$.
\end{theorem}

\begin{proof}
$U = (13/4, 15/4)$ is an open subset of $(3,\infty)$, and it is nonempty because $7/2 \in U$. If $V$ were
dense in $(3,\infty)$, then $U$ would contain some $r \in V$, i.e.\ $r = \repdim A$ for some
finite-dimensional algebra $A$. By Lemma~\ref{lem:int}, $r$ is an integer, and by Lemma~\ref{lem:gap},
$r \notin U$. This is a contradiction.
\end{proof}

\begin{corollary}
Conjecture 00000004042 is false: for every proposition $P$, in particular for clause (2) however it is
formalized, $\neg(\text{(1)} \wedge P)$.
\end{corollary}

\section{Lean 4 formalization}\label{sec:lean}

The directory \texttt{lean/} contains a Lean~4 project (Lean \texttt{v4.33.1}, Mathlib \texttt{v4.33.1}).
\texttt{Conjecture4042.lean} imports the single source file \texttt{Conjecture4042/Basic.lean}
(228 lines, namespace \texttt{Conjecture4042}, one universe variable \texttt{u}), which contains the
following.

\paragraph{Definitions.}
\begin{itemize}
  \item \texttt{globalDimension R} $:= \sup_{X}$ \texttt{projectiveDimension X}, the supremum
        (Lean \texttt{iSup}) over all \texttt{X : ModuleCat.\{u\} R}, valued in
        \texttt{WithBot $\mathbb{N}_\infty$} (Definition~\ref{def:gldim}).
  \item \texttt{DualRegular k A}, the type of $k$-linear maps $A \to k$, i.e.\ $D(A_A)=\Hom_k(A,k)$, with the left $A$-module
        structure $(a\cdot f)(x) = f(xa)$, defined as \texttt{f.comp (LinearMap.mulRight k a)}.
  \item \texttt{IsSummandOfPower X M}: there are $n$ and $A$-linear maps
        \texttt{i : X $\to$ (Fin n $\to$ M)}, \texttt{r : (Fin n $\to$ M) $\to$ X} with
        \texttt{r $\circ$ i = id}, i.e.\ $X \in \add M$.
  \item \texttt{IsGeneratorCogenerator k A M}: \texttt{Module.Finite A M}, \texttt{A $\in$ add M} and
        \texttt{D(A) $\in$ add M}.
  \item \texttt{repdim k A} $:= \inf_{M}$ \texttt{globalDimension (Module.End A M)}, the infimum
        (Lean \texttt{iInf}) over all \texttt{M : ModuleCat.\{u\} A} with
        \texttt{IsGeneratorCogenerator k A M} (Definition~\ref{def:repdim}).
  \item \texttt{toEReal : WithBot $\mathbb{N}_\infty$ $\to$ EReal}, the order embedding
        $\bot\mapsto\bot$, $n\mapsto n$, $\top\mapsto\top$, used only to compare values of
        \texttt{repdim} with real numbers.
  \item \texttt{IsRealizable v}: there are \texttt{k A : Type u} with \texttt{[Field k] [Ring A]
        [Algebra k A] [FiniteDimensional k A]} and \texttt{toEReal (repdim k A) = v}.
  \item \texttt{DenseAbove3 := Dense \{x : Set.Ioi (3:$\mathbb{R}$) | IsRealizable ((x:$\mathbb{R}$) :
        EReal)\}}: clause (1), density in the subspace $(3,\infty)$ of $\mathbb{R}$.
  \item \texttt{DerivedEquivalent A B}: there is an equivalence
        \texttt{e : DerivedCategory (ModuleCat A) $\simeq$ DerivedCategory (ModuleCat B)} of Mathlib's
        derived categories whose functor commutes with the shift
        (\texttt{e.functor.CommShift $\mathbb{Z}$}) and is triangulated
        (\texttt{e.functor.IsTriangulated}). Mathlib provides the derived category of all modules. We
        use it because it is the derived category available in Mathlib for a noncommutative ring. This
        choice is not load-bearing; see Section~\ref{sec:reading}.
  \item \texttt{ClosedUnderDerivedEquivalence}: for all finite-dimensional algebras $A$ over $k$ and $B$
        over $k'$, \texttt{repdim k A = 3 $\to$ DerivedEquivalent A B $\to$ repdim k' B = 3}
        (clause (2)).
\end{itemize}

\paragraph{Theorems.}
\begin{itemize}
  \item \texttt{exists\_nat\_of\_toEReal\_eq\_coe}: if \texttt{toEReal x = r} for a real $r$, then
        \texttt{x = m} and \texttt{r = m} for some \texttt{m : $\mathbb{N}$}.
  \item \texttt{repdim\_eq\_nat\_of\_toEReal\_eq}, \texttt{eq\_nat\_of\_isRealizable}:
        Lemma~\ref{lem:int} (every real realizable value of \texttt{repdim} is a natural number).
  \item \texttt{no\_nat\_in\_Ioo}: Lemma~\ref{lem:gap}.
  \item \texttt{not\_isRealizable\_of\_mem\_Ioo}: no realizable value lies in $(13/4,15/4)$.
  \item \texttt{not\_denseAbove3 : $\neg$ DenseAbove3} (Theorem~\ref{thm:main}), via
        \texttt{dense\_iff\_inter\_open} applied to the open set
        $\{x \in (3,\infty) : 13/4 < x < 15/4\}$, which contains $7/2$.
  \item \texttt{conjecture\_00000004042\_false\_of\_any\_clause2 (P : Prop) : $\neg$ (DenseAbove3 $\wedge$ P)}.
  \item The main theorem
\[
  \texttt{conjecture\_00000004042\_false} :\ \neg\,(\texttt{DenseAbove3} \wedge
  \texttt{ClosedUnderDerivedEquivalence}).
\]
\end{itemize}

\paragraph{Axiom audit.} There is no \texttt{sorry}, no \texttt{native\_decide} and no added axiom.
\texttt{\#print axioms conjecture\_00000004042\_false} reports
\begin{verbatim}
'Conjecture4042.conjecture_00000004042_false' depends on axioms:
  [propext, Classical.choice, Quot.sound]
\end{verbatim}
and the same holds for \texttt{conjecture\_00000004042\_false\_of\_any\_clause2}. The file also
elaborates with \texttt{-DwarningAsError=true}. To check:
\begin{verbatim}
cd lean
lake exe cache get
lake build
\end{verbatim}

\end{document}
